% Chapter 2 - State of the Art
% Rebuilt for clarity: each section asks one question and ends with a "For InvestiGator" takeaway.
% All citations, tables and figures preserved. EN-GB.

\chapter{State of the Art}
\label{chap:sota}

This chapter has one job: to justify every design choice in InvestiGator against the alternatives. Each section
asks what a field offers and ends with what InvestiGator takes from it. The order is: how retail investors
behave; anomaly detection; explainable AI; trust in automation; financial text representation; information
retrieval; event studies; and case-based reasoning. The recurring finding is simple. The building blocks
are mature, but an integrated, explainable system for the retail investor is still rare.

\section{Retail Investing and Information Overload}
\label{sec:sota_retail}

\emph{How do retail investors behave, and what should an alert give them?}

Behavioural finance shows that real investors depart systematically from the rational ideal
\autocite{barberis2003behavioral}. One root cause is \emph{loss aversion}: we feel losses more than equal
gains \autocite{kahneman1979prospect}, which makes investors especially sensitive to bad news, and a
timely, well-explained alert valuable.

The strongest pattern is \emph{attention}. \textcite{barber2008glitters} show that individuals are net
buyers of attention-grabbing stocks (those in the news, with unusual volume, or extreme one-day returns):
facing thousands of choices, they fall back on whatever caught their eye. Attention can even be measured
through search volume \autocite{da2011attention}. This is decisive for the design, because a sharp price
move and a salient headline are exactly InvestiGator's two triggers, so an alert that adds context lands where
attention-driven decisions are made.

News also moves markets, and its text can be measured. \textcite{tetlock2007media} builds a daily
pessimism index from a newspaper column and links it to price pressure and trading volume, early
evidence that textual signals relate to market outcomes, which is the premise of InvestiGator's news
engine. This
audience is large and active too: even through the March~2020 crash, retail investors on one
commission-free platform added to their holdings rather than panicking \autocite{welch2022robinhood},
consistent with the broad ownership and market scale reported in Chapter~\ref{chap:introduction}
\autocite{gallup2025stock, sifma2025factbook}.

\textbf{For InvestiGator:} retail investors are attention-constrained and news-driven, so the system does not
predict the market; it lowers the cost of \emph{interpreting} the events that already grab attention, with
a transparent explanation and real historical evidence.

\section{Anomaly Detection in Financial Markets}
\label{sec:sota_anomaly}

\emph{How do you tell an abnormal price move from ordinary noise?}

InvestiGator's first trigger is an anomaly detector, so it sits in the mature field of anomaly detection. The
standard taxonomy \autocite{chandola2009anomaly} separates \emph{point}, \emph{contextual} and
\emph{collective} anomalies, and groups methods into statistical, distance/density, and machine-learning
families (Figure~\ref{fig:sota_anomaly_taxonomy}). A sharp single-day return is a contextual point
anomaly: abnormal relative to the asset's own recent behaviour.

\begin{figure}[ht]
\centering
\begin{tikzpicture}[
  font=\small,
  node distance=4mm,
  root/.style={draw, rounded corners, fill=black!5, align=center, text width=3.2cm, minimum height=0.8cm},
  fam/.style={draw, rounded corners, align=center, text width=3.0cm, minimum height=1.4cm},
  >={Stealth[length=2mm]}
]
  \node[root] (r) {Anomaly detection methods};
  \node[fam, below left=10mm and 1.2cm of r] (stat) {\textbf{Statistical}\\ \footnotesize z-score, control charts; transparent, assume a distribution};
  \node[fam, below=10mm of r] (ml) {\textbf{Distance / density}\\ \footnotesize Isolation Forest, LOF; flexible, less interpretable};
  \node[fam, below right=10mm and 1.2cm of r] (dl) {\textbf{Deep learning}\\ \footnotesize autoencoders, deep one-class; powerful, opaque, data-hungry};
  \draw[->] (r) -- (stat);
  \draw[->] (r) -- (ml);
  \draw[->] (r) -- (dl);
\end{tikzpicture}
\caption[Taxonomy of anomaly-detection methods]{Families of anomaly-detection methods relevant to this
work, following \textcite{chandola2009anomaly} and \textcite{pang2021deep}. Interpretability decreases,
and modelling capacity increases, from left to right.}
\label{fig:sota_anomaly_taxonomy}
\end{figure}

The simplest detector is a rolling \emph{z}-score: flag a return that deviates far from its recent mean and
standard deviation. It is attractive because it is transparent and easy to justify. Its one assumption, a
locally stable distribution, is handled by the rolling window, which tracks \emph{volatility clustering},
the well-known tendency of calm and turbulent periods to arrive in runs
\autocite{engle1982arch, bollerslev1986garch}. The rolling standard deviation is, in effect, a simpler
non-parametric version of that idea: it adapts to the current regime without fitting a model the user
could not scrutinise. Between the two sits the exponentially weighted moving average of squared returns
(the RiskMetrics estimator), which weights the recent past more heavily than the distant past but still
estimates no parameters per instrument. A full GARCH fit is declined here under defensible simplicity,
but the claim is not left as opinion: the EWMA variant, the first step on the road to GARCH, is put
through the same evaluation protocol in Chapter~\ref{chap:casestudies}, so the trade-off is measured
rather than asserted.

More expressive detectors trade interpretability for capacity. Isolation Forest isolates outliers by
random partitioning and scales to high dimensions \autocite{liu2008isolation}; the Local Outlier Factor
scores how isolated a point is within its neighbourhood \autocite{breunig2000lof}; deep detectors learn a
model of normality \autocite{pang2021deep}. All are powerful but harder to read off, and their edge shows
most in high-dimensional data, not a single return series. Rather than dismissing them on those grounds
alone, Chapter~\ref{chap:casestudies} evaluates both Isolation Forest and the Local Outlier Factor under
the same causal protocol as the statistical rule, so the comparison is empirical. In finance, detection
is anyway dominated by \emph{unsupervised} methods, because labelled anomalies are scarce
\autocite{ahmed2016financial}, a constraint that also shapes how InvestiGator is evaluated
(Chapter~\ref{chap:casestudies}). Table~\ref{tab:sota_anomaly} compares the families.

\begin{table}[ht]
\caption{Anomaly-detection approaches relevant to this work.}
\label{tab:sota_anomaly}
\centering
\begin{tabular}{p{0.24\textwidth} p{0.30\textwidth} p{0.38\textwidth}}
\toprule
\tabhead{Approach} & \tabhead{How it works} & \tabhead{Strengths / limitations} \\
\midrule
Statistical \emph{z}-score (rolling) \autocite{chandola2009anomaly} & Flags returns that deviate from a rolling mean and standard deviation & Transparent and easy to explain; assumes a locally stable distribution \\
Isolation Forest \autocite{liu2008isolation} & Isolates points via random partitioning (ensemble) & Scalable, handles high dimensionality; score is not directly interpretable \\
Local Outlier Factor \autocite{breunig2000lof} & Scores isolation within a local density neighbourhood & Captures local structure; sensitive to neighbourhood size, score hard to read \\
Deep detectors \autocite{pang2021deep} & Learn a model of normality (e.g.\ autoencoders) & High capacity for complex data; opaque, data-hungry, hard to justify \\
\bottomrule
\end{tabular}
\end{table}

\textbf{For InvestiGator:} the signal is one low-dimensional return series that must be explained to a
non-expert, so the transparent \emph{z}-score wins; the more expressive families are noted as alternatives
whose opacity the problem does not justify.

\section{Explainable Artificial Intelligence}
\label{sec:sota_xai}

\emph{What kind of explanation should the system give?}

Explainable AI (XAI) exists because the strongest models rarely show their reasoning. Surveys map the field
and its central split \autocite{adadi2018peeking, arrieta2020xai}: \emph{transparent} models, interpretable
in themselves, versus \emph{post-hoc} methods that explain a trained black box, which can be organised by
what they explain \autocite{guidotti2018survey} (Figure~\ref{fig:sota_xai_taxonomy}). A useful caution:
``interpretability'' is not one thing \autocite{lipton2018mythos}, so a system should say which kind it
offers. InvestiGator offers transparency of its decision logic and traceable evidence, not a claim of full
interpretability.

\begin{figure}[ht]
\centering
\begin{tikzpicture}[
  font=\small,
  node distance=4mm,
  root/.style={draw, rounded corners, fill=black!5, align=center, text width=3.0cm, minimum height=0.8cm},
  br/.style={draw, rounded corners, align=center, text width=3.4cm, minimum height=1.5cm},
  >={Stealth[length=2mm]}
]
  \node[root] (r) {Explainability};
  \node[br, below left=10mm and 0.6cm of r] (tr) {\textbf{Transparent design}\\ \footnotesize inherently interpretable logic (rules, linear models, retrieved evidence)};
  \node[br, below right=10mm and 0.6cm of r] (ph) {\textbf{Post-hoc}\\ \footnotesize explain a trained black box (LIME, SHAP, counterfactuals)};
  \draw[->] (r) -- (tr);
  \draw[->] (r) -- (ph);
\end{tikzpicture}
\caption[Taxonomy of explainability approaches]{The two broad routes to explainability
\autocite{arrieta2020xai, guidotti2018survey}: building models that are transparent by design, or
explaining a black box after the fact. This work follows the former.}
\label{fig:sota_xai_taxonomy}
\end{figure}

Two post-hoc techniques are widely used: LIME fits a simple local surrogate around one prediction
\autocite{ribeiro2016lime}, and SHAP attributes a prediction to its features using Shapley values
\autocite{lundberg2017shap}. Both are valuable when one is committed to a black box. But for high-stakes
decisions, \textcite{rudin2019stop} argues for dropping that commitment and using models that are
interpretable in the first place, since a post-hoc explanation is itself an approximation that can mislead.
InvestiGator follows this: it explains primarily by transparent design (every quantity exposed, real precedents
shown as evidence), with feature attribution as an optional extra.

Good explanations are also a human matter. People want \emph{contrastive}, selective reasons, not an
exhaustive dump \autocite{miller2019explanation}, which is why each alert shows a few relevant precedents
and the exact quantities behind the decision. Explanations must also be evaluated, not assumed:
\textcite{doshivelez2017rigorous} argue that interpretability claims need rigorous evaluation rather
than assertion, and set out how such evaluation can be structured. The criterion this work adopts is
\emph{fidelity}, how well an explanation reflects the real computation, because it is the property
InvestiGator can guarantee (Chapter~\ref{chap:casestudies}). Table~\ref{tab:sota_xai} summarises the
options.

\begin{table}[ht]
\caption{Explainability approaches considered.}
\label{tab:sota_xai}
\centering
\begin{tabular}{p{0.24\textwidth} p{0.30\textwidth} p{0.38\textwidth}}
\toprule
\tabhead{Method} & \tabhead{Type} & \tabhead{Strengths / limitations} \\
\midrule
LIME \autocite{ribeiro2016lime} & Local surrogate, post-hoc & Intuitive, model-agnostic; explanations can be unstable, only locally faithful \\
SHAP \autocite{lundberg2017shap} & Shapley-value attribution, post-hoc & Consistent, theoretically grounded; computationally costly \\
Transparent design \autocite{rudin2019stop, arrieta2020xai} & Inherently interpretable logic & No approximation, faithful by construction; constrains model choice \\
\bottomrule
\end{tabular}
\end{table}

\textbf{For InvestiGator:} pursue explainability by transparent design and judge it by fidelity; treat post-hoc
attribution as a complement, not the foundation.

\section{Trust and Appropriate Reliance in Decision Support}
\label{sec:sota_trust}

\emph{When should the investor actually trust an alert?}

An explanation only helps if it produces \emph{appropriate} reliance: not dismissing a sound warning, not
acting blindly on a fallible one. The human-factors view \autocite{lee2004trust} calls the goal
\emph{calibrated} trust, reliance matched to the system's real competence, and names two failure modes,
\emph{disuse} (ignoring a good aid) and \emph{misuse} (over-trusting a flawed one). With real money, both
are costly.

Explanations are how trust gets calibrated, but they are not magic. \textcite{bansal2021whole} show that
explanations can even cause \emph{over-reliance}, with people accepting confident explanations of wrong
answers. InvestiGator's response is concrete: explain with \emph{verifiable evidence} (the actual $z$-score and
its parameters; real precedents with their measured impact), and disclaim prediction on every alert.

\textbf{For InvestiGator:} transparency is not decoration but the mechanism for appropriate reliance, so the
system shows checkable evidence rather than a persuasive story.

\section{Financial NLP and Text Representation}
\label{sec:sota_finnlp}

\emph{How should the system represent the meaning of a headline?}

To compare headlines, text must become numbers. The financial-NLP literature falls into three
generations: lexicons, word embeddings, and contextual neural models. Surveys of textual analysis in
finance up to the mid-2010s \autocite{kearney2014textual} cover the first two; the third arrived
afterwards, and is where this work takes its representation.

Lexicons came first. Finance-specific word lists fix the fact that general sentiment dictionaries misread
ordinary financial words such as ``liability'' or ``tax'' \autocite{loughran2011liability}; they are fully
transparent but bounded by their vocabulary. Word embeddings then placed words in a vector space where
proximity means similar use: word2vec from local context \autocite{mikolov2013word2vec}, GloVe from global
co-occurrence \autocite{pennington2014glove}. Both give each word a single vector, blind to context.

The Transformer \autocite{vaswani2017attention} removed that limit, and BERT \autocite{devlin2019bert}
produced deeply \emph{contextual} representations. Comparing two texts with BERT directly is costly, though;
Sentence-BERT \autocite{reimers2019sbert} fixes this by giving each sentence one fixed-length vector
comparable by a cheap cosine, exactly what precedent retrieval needs. At very large scale, exact search can
be swapped for an approximate index \autocite{johnson2021faiss} without changing the representation (future
work). Two heavier options also exist: domain-tuned FinBERT models
\autocite{araci2019finbert, yang2020finbert} and large financial LLMs such as BloombergGPT
\autocite{wu2023bloomberggpt}, both powerful but costlier, less transparent, and outside this work's free,
reproducible constraint. Table~\ref{tab:sota_nlp} compares them.

\begin{table}[ht]
\caption{Text-representation approaches for financial news.}
\label{tab:sota_nlp}
\centering
\begin{tabular}{p{0.22\textwidth} p{0.32\textwidth} p{0.38\textwidth}}
\toprule
\tabhead{Approach} & \tabhead{Representation} & \tabhead{Strengths / limitations} \\
\midrule
Finance lexicons \autocite{loughran2011liability} & Curated word lists / tone & Fully transparent; bounded by vocabulary \\
word2vec / GloVe \autocite{mikolov2013word2vec, pennington2014glove} & Static word vectors & Efficient; one vector per word, no context \\
BERT \autocite{vaswani2017attention, devlin2019bert} & Contextual token embeddings & Strong understanding; pairwise similarity is costly \\
Sentence-BERT \autocite{reimers2019sbert} & Sentence embeddings (siamese) & Efficient cosine similarity; the choice in this work \\
FinBERT \autocite{araci2019finbert, yang2020finbert} & Domain-adapted language model & Finance-tuned; oriented to sentiment / communications \\
Financial LLMs \autocite{wu2023bloomberggpt} & Large generative model & Very capable; costly, opaque, outside free-tier scope \\
\bottomrule
\end{tabular}
\end{table}

\begin{figure}[ht]
\centering
\begin{tikzpicture}[
  font=\footnotesize, >={Stealth[length=2.4mm]},
  gen/.style={draw, rounded corners, align=center, text width=2.7cm, inner sep=4pt,
              minimum height=2.5cm, anchor=south},
]
  \node[gen, fill=blue!4, draw=blue!40] (g1) at (0,0)
    {\textbf{Lexicons}\\[2pt] curated word lists\\[4pt] \emph{transparent, but}\\ \emph{vocabulary-bound}};
  \node[gen, fill=teal!6, draw=teal!45] (g2) at (3.3,0)
    {\textbf{word2vec / GloVe}\\[2pt] static word vectors\\[4pt] \emph{one vector per}\\ \emph{word, no context}};
  \node[gen, fill=orange!8, draw=orange!55] (g3) at (6.6,0)
    {\textbf{BERT}\\[2pt] contextual token\\ embeddings\\[4pt] \emph{deep context, but}\\ \emph{costly to compare}};
  \node[gen, fill=green!12, draw=green!55!black, line width=1pt] (g4) at (9.9,0)
    {\textbf{Sentence-BERT} $\star$\\[2pt] one vector per\\ sentence\\[4pt] \emph{cheap cosine,}\\ \emph{the choice here}};
  \draw[->, thick, black!55] (-1.45,-0.7) -- (11.35,-0.7);
  \node[below, font=\small, black!65] at (4.95,-0.75)
    {increasingly contextual meaning; Sentence-BERT keeps comparison cheap};
\end{tikzpicture}
\caption[Three generations of text representation]{Three generations of turning financial text into
numbers. Finance lexicons are transparent but bounded by their vocabulary; static word vectors add
similarity of use but no context; contextual models (BERT) capture meaning yet are costly to compare
pairwise; Sentence-BERT gives each headline a single vector comparable by a cheap cosine, exactly what
precedent retrieval needs, and is the representation chosen here. Details in Table~\ref{tab:sota_nlp}.}
\label{fig:sota_nlp_ladder}
\end{figure}

\textbf{For InvestiGator:} Sentence-BERT is the sweet spot (Figure~\ref{fig:sota_nlp_ladder}),
contextual meaning beyond lexicons and static vectors, yet cheap, comparable and free to run. One honesty note on the comparison: the evaluation in
Chapter~\ref{chap:casestudies} measures Sentence-BERT against a \emph{lexical} baseline, which brackets
the question that matters (does meaning beat surface words?); the embedder choice itself was also
benchmarked on the same protocol. A domain-tuned FinBERT, mean-pooled into sentence embeddings, reaches
only a precision@5 of $0.420$, below MiniLM's $0.514$ and consistent with its sentiment orientation,
while current general encoders (E5-small, BGE-small) tie MiniLM (around $0.51$) rather than beating it.
The small, free, general encoder is therefore the sweet spot by measurement, not merely by argument;
only static word vectors, strictly dominated for sentence comparison, are left to argument.

\section{Information Retrieval and Ranking Evaluation}
\label{sec:sota_ir}

\emph{Once headlines are vectors, how do you rank past news, and how do you tell if the ranking is good?}

Ranking a corpus by relevance to a query is the defining problem of information retrieval (IR). The
vector-space model \autocite{salton1975vsm} represents documents and queries as vectors ranked by cosine
similarity; InvestiGator's embedding retrieval is the same idea with richer vectors. Classical \emph{lexical}
ranking weights informative terms, from TF-IDF to BM25 \autocite{robertson2009bm25}, still a strong
baseline. Its known weakness is \emph{vocabulary mismatch}: a relevant document that uses different words
scores as dissimilar. Semantic embeddings overcome exactly that, which is why InvestiGator compares against a
lexical baseline (Chapter~\ref{chap:casestudies}).

Because retrieval returns a ranked list, it is judged with rank-aware measures \autocite{manning2008ir}; the
most direct is \emph{precision@$k$}, the fraction of the top $k$ results that are relevant. InvestiGator adopts
it because an alert can show only a handful of precedents, so what matters is whether the top few are
relevant, reported against random, recency and lexical baselines.

\textbf{For InvestiGator:} retrieval reuses a mature representation-then-rank pipeline and the standard
precision@$k$ metric; the contribution is the application and the pairing with measured impact, not a new
algorithm.

\section{Event Studies and News--Market Impact}
\label{sec:sota_eventstudy}

\emph{Having found analogous news, how do you measure what it did to the price, honestly?}

The tool is the \emph{event study}, introduced by \textcite{fama1969adjustment} and standardised for daily
returns by \textcite{brown1985daily}, who set the practice of measuring returns over short post-event
windows, the basis for InvestiGator's $+1$, $+3$ and $+5$-day horizons. The impact is measured strictly
\emph{after} the event, so it describes an outcome, not a forecast.

Why not just predict the price? Because public news is absorbed into prices almost at once, the
efficient-market hypothesis \autocite{fama1970efficient}, so forecasting returns from public news is very
hard by construction. InvestiGator therefore does not compete with market efficiency; it measures and explains
the reactions that \emph{already} followed comparable news, which is both more honest and more defensible
for a non-expert.

Relating news \emph{text} to impact is an active area: \textcite{tetlock2007media} is an early example, and
a more recent strand tries to \emph{predict} moves from news \autocite{xing2018nlffsurvey, ding2015deep}.
InvestiGator deliberately stops short of prediction. Doing this at scale needs news aligned to prices, which is
the contribution of the FNSPID dataset \autocite{dong2024fnspid}, used to build the knowledge base; its
coverage and period limits are noted in the data card of Chapter~\ref{chap:methods}.

\textbf{For InvestiGator:} pairing semantic retrieval with event-study impact over FNSPID is the methodological
core, evidence about the past, never a forecast.

\section{Case-Based Reasoning and Precedent Retrieval}
\label{sec:sota_cbr}

\emph{What is the engine, in one familiar idea?}

It is \emph{case-based reasoning} (CBR). In the classic account \autocite{aamodt1994cbr}, a CBR system
\emph{retrieves} similar past cases, \emph{reuses} their outcomes, and optionally \emph{revises} and
\emph{retains} them. InvestiGator is the retrieve-and-reuse core: each past headline plus its measured impact is
a \emph{case}; a new headline is the query; cosine similarity retrieves; the precedents' impact is reused as
evidence. It stops before \emph{revise}, so it never turns evidence into a prediction.

\textbf{For InvestiGator:} naming the engine as CBR shows it is a recognised paradigm, not ad hoc, and its
``this resembles these past cases'' explanations are naturally intelligible to a non-expert.

\section{Learned Alert Triage: Supervised Materiality Classification}
\label{sec:sota_triage}

\emph{When dozens of headlines arrive in a day, which ones deserve an alert?}

Retrieval answers ``what does this resemble?'', and the event study answers ``what followed similar
news?''. Neither ranks the incoming stream. That is a \emph{triage} problem, and it can be posed as
supervised classification without breaking the no-forecast rule. The model estimates the probability
that an \emph{abnormal move of unusual size, in either direction}, follows a news item; the label is
produced by the system's own event-study measurement, the market-adjusted return of
\textcite{brown1985daily} against an index proxy. The target is \emph{materiality}, whether the market
reacted at all. Direction and magnitude of the price path are never predicted, so the task stays on
the defensible side of the efficient-market line drawn earlier in this chapter.

Two model families cover the interpretability--capacity trade-off. Logistic regression is the
transparent choice: with standardised inputs, its log-odds are an exact additive sum of per-feature
contributions. That makes it the kind of inherently interpretable model
\textcite{rudin2019stop} argues should be preferred when explanations matter. Gradient-boosted decision
ensembles \autocite{friedman2001gbm} are the stronger nonlinear learner and serve as the ceiling the
interpretable model is compared against. For an end user, the score is only honest if it is a
\emph{probability}: raw classifier scores are typically miscalibrated, and post-hoc calibration, such as
fitting a sigmoid on held-out validation data (Platt scaling), corrects this
\autocite{niculescu2005calibration}. Under class imbalance, evaluation uses the area under the
precision--recall curve rather than accuracy. A budget-shaped precision completes the picture (how
many of the $N$ alerts a user can absorb per day were material), because that budget is exactly the
alert-fatigue cost the triage exists to reduce.

\textbf{For InvestiGator:} triage adds a \emph{trained} severity score on top of the transparent pipeline,
ranked and gated per day, explained by its additive contributions, calibrated on validation data, and
always compared against a volatility-only baseline before it is allowed to matter.

\section{Uncertainty and Drift in a Deployed Model}
\label{sec:sota_uncertainty}

\emph{A calibrated score is still not a guarantee, and a model does not stay still. What does the
literature offer on both counts?}

Calibration, as used for the triage score above, is an \emph{aggregate} property: among items scored
near a given probability, roughly that fraction should turn out positive. It says nothing about any
individual item, and it offers no recourse once the model is applied outside the conditions in which
it was fitted. \emph{Conformal prediction} narrows this gap. It returns a set for each item rather
than a point estimate, although its guarantee remains \emph{marginal}, averaged over cases, rather
than conditional on any one of them. The foundational treatment
\autocite{vovk2005algorithmic} shows that a predictor can be wrapped so that it returns a
\emph{set} of labels carrying a guarantee that is distribution-free and valid in finite samples: for
a chosen error rate $\alpha$, the set contains the true label in at least $1-\alpha$ of cases. The
guarantee assumes neither normality nor a correctly specified model, and requires only that
calibration and test data be exchangeable. The split (or inductive) variant, which calibrates on a
held-out block and so costs one extra pass rather than a retraining loop,
is set out with its finite-sample quantile correction by \textcite{angelopoulos2023conformal}.

The appeal for a system built on refusing to overclaim is direct. In a binary decision the returned
set may contain \emph{both} labels, which is an explicit and quantified ``I do not know'' rather than
a probability near one half that still forces a choice. This is a different kind of honesty from
calibration, and the two are complementary rather than alternatives.

The exchangeability assumption then raises the second question. Models are trained on a past and
deployed into a future, and the relationship between inputs and target can change underneath them.
The standard treatment of this in machine learning is the concept-drift literature, whose survey
\autocite{gama2014survey} sets out the taxonomy (sudden, gradual, incremental and recurring change),
the detection strategies and the adaptation methods. It is a general treatment rather than a
financial one, so it supplies the framing and the vocabulary while any statement about how much a
particular corpus has drifted must come from measuring that corpus.

Both concerns belong to a wider point about systems rather than models. \textcite{sculley2015debt}
argue, from experience with production systems, that machine-learning components accumulate a
distinctive and largely hidden maintenance cost: entanglement, where changing anything changes
everything; unstable and undeclared data dependencies; and configuration debt. Their account is a
position paper drawn from practice rather than an empirical study, and it is cited here for that
qualitative claim alone. Its relevance is that a dissertation prototype which runs continuously,
as this one does, inherits the same category of problem in miniature.

\textbf{For InvestiGator:} these three strands convert two of the system's standing caveats into
measurable quantities. The claim that a probability is not a promise becomes a coverage experiment
(Section~\ref{sec:cs_conformal}); the claim that the training corpus predates deployment becomes a
drift measurement with conventional interpretation bands (Section~\ref{sec:cs_drift}). Neither
changes what the deployed system does, and in both cases that restraint is itself the reported
result.

\section{Existing Tools for the Retail Investor}
\label{sec:sota_tools}

\emph{What can a retail investor already use, and what is missing?}

Three kinds of tool are common today, and each leaves a gap. \emph{Price alerts}, in every brokerage app,
fire when a price or percentage threshold is crossed; they ride the same attention behaviour InvestiGator
targets \autocite{barber2008glitters}, but a fixed threshold ignores each stock's own volatility (firing
constantly on volatile names, rarely on calm ones) and says only \emph{that} a level was crossed, never
\emph{why}. \emph{News and sentiment apps} aggregate headlines and attach a sentiment score, building on the
link between tone and markets \autocite{tetlock2007media} and on finance lexicons
\autocite{loughran2011liability}; but a single polarity is thin, often proprietary, and rarely connects a
headline to concrete historical precedents and their outcomes. \emph{Robo-advisors}, the most studied
retail-fintech tool, genuinely help by improving diversification and curbing biases
\autocite{dacunto2019robo, cardillo2024robo}, but they \emph{act for} the investor with proprietary logic,
the opposite of an explanation-first tool that leaves the decision to the user. Table~\ref{tab:sota_tools}
positions all three against InvestiGator.

\begin{table}[ht]
\caption{Existing retail tools versus the system proposed in this work.}
\label{tab:sota_tools}
\centering
\begin{tabular}{p{0.24\textwidth} p{0.18\textwidth} p{0.16\textwidth} p{0.16\textwidth} p{0.11\textwidth}}
\toprule
\tabhead{Tool} & \tabhead{What it does} & \tabhead{Explains why?} & \tabhead{Shows precedents?} & \tabhead{Acts for user?} \\
\midrule
Brokerage price alert & Fires on a fixed price / \% threshold & No & No & No \\
News / sentiment app \autocite{tetlock2007media} & Aggregates news, scores sentiment & Partly (opaque score) & No & No \\
Robo-advisor \autocite{dacunto2019robo} & Allocates and rebalances a portfolio & Rarely (proprietary) & No & Yes \\
\textbf{This work (InvestiGator)} & Explains events, shows evidence & Yes, by design & Yes (measured impact) & No \\
\bottomrule
\end{tabular}
\end{table}

A fourth option has become common since large language models reached consumer products: asking a
general assistant why a stock moved. These models are demonstrably capable on financial text, and
domain-specific variants exist \autocite{wu2023bloomberggpt}. The difficulty is not fluency but
grounding. A general assistant has no access to the security's own recent volatility, so it cannot
say whether a move was unusual \emph{for that stock}; it does not estimate the security's sensitivity
to its market and sector, so it cannot separate company from market; and its account of comparable
past events is recalled rather than retrieved from an indexed corpus with measured outcomes. The
answer reads well and cannot be checked, which is precisely the interpretability failure the
explainability literature warns about \autocite{lipton2018mythos, rudin2019stop}.

Since 2025 that fourth option has stopped being a workaround and become a shipped feature, which
makes it the closest competitor to this work's central claim and worth naming rather than describing
as a category. Robinhood introduced \emph{Cortex} in March~2025 with a feature whose stated purpose is
to ``answer the age-old question of, `Why is this stock going up or down today?'\,'' by generating
summaries from news, analyst reports, technical signals and the broker's own platform data
\autocite{robinhood2025cortex}. Google rebuilt Google Finance around AI and shipped
``key moments'' that ``explain why a stock moved'', annotated onto the price chart
\autocite{google2026finance}. Both were observed on 2026-08-07; the wording quoted is each vendor's
own.

The overlap with this work is real and should be stated plainly: these products answer the same
question for a far larger audience than a dissertation ever will, and they answer it in plain
language, which was one of the goals here. The difference is what stands behind the answer, and it
is a difference of \emph{kind}. A generated summary is a claim; this work's output is a claim with
its evidence attached. Neither product states whether the move was unusual \emph{for that security}
against its own recent volatility, neither separates the company's contribution from its market and
sector, and neither retrieves comparable past events with outcomes measured after the fact. Nor is
either verifiable by the reader: Robinhood's own disclosure is that ``there is no guarantee that AI
will improve investing performance, mitigate risk, or reduce losses''
\autocite{robinhood2025cortex}, which is a statement about outcomes and not about whether any
individual explanation is right. That is the gap this work occupies, and it is a narrow one honestly
described: not a better answer, but an answer a reader can check.

Table~\ref{tab:sota_three_questions} therefore scores the same tools against the three questions
posed in Section~\ref{sec:intro_problem}, which is a stricter test than a feature list because each
question has a checkable answer.

\begin{table}[ht]
\caption{Existing tools scored against the three questions a retail investor asks
(Section~\ref{sec:intro_problem}).}
\label{tab:sota_three_questions}
\centering
\begin{tabular}{p{0.26\textwidth} p{0.19\textwidth} p{0.19\textwidth} p{0.22\textwidth}}
\toprule
\tabhead{Tool} & \tabhead{Q1 Unusual for this stock?} & \tabhead{Q2 Company or market?} & \tabhead{Q3 Happened before, with outcomes?} \\
\midrule
Brokerage price alert & No (fixed threshold) & No & No \\
News / sentiment app & No & No & No (polarity only) \\
Robo-advisor \autocite{dacunto2019robo} & Not surfaced & Not surfaced & No \\
General LLM assistant \autocite{wu2023bloomberggpt} & No (no volatility baseline) & No (no beta estimate) & Recalled, not retrieved \\
AI ``why did it move?'' feature (Robinhood Cortex \autocite{robinhood2025cortex}; Google Finance key moments \autocite{google2026finance}) & Not stated & Not stated & Generated, not retrieved \\
Professional terminal & Yes & Yes & Yes (paid) \\
\textbf{This work (InvestiGator)} & Yes (rolling z-score) & Yes (rolling beta) & Yes (retrieved, measured) \\
\bottomrule
\end{tabular}
\end{table}

\textbf{For InvestiGator:} each existing free tool covers one corner, alerting, news, or allocation,
and none combines volatility-aware detection, market-versus-company attribution and concrete
historical evidence in a single tool that informs rather than decides. Professional terminals do
combine them, at a price the target user does not pay. That intersection, all three answers at no
cost and with the reasoning exposed, is the niche this work fills.

\section{Comparative Analysis and Positioning}
\label{sec:sota_positioning}

\emph{So what does InvestiGator choose, and why?}

Table~\ref{tab:sota_positioning} gathers the choices against their alternatives. The rule throughout is
\emph{defensible simplicity} \autocite{rudin2019stop}: when two options perform comparably, prefer the more
transparent one, because the system must be understood and acted upon by a non-expert.

\begin{table}[ht]
\caption{Component choices in this work versus alternatives.}
\label{tab:sota_positioning}
\centering
\begin{tabular}{p{0.20\textwidth} p{0.26\textwidth} p{0.22\textwidth} p{0.22\textwidth}}
\toprule
\tabhead{Component} & \tabhead{Chosen} & \tabhead{Alternative} & \tabhead{Why} \\
\midrule
Anomaly detection & \emph{z}-score (rolling) & Isolation Forest, deep detectors & Transparency; low-dimensional signal \\
News retrieval & Sentence-BERT + cosine & Lexicons, word2vec, LLMs & Contextual yet cheap and reproducible \\
Impact measurement & Event-study windows & --- & Standard, reproducible, no forecasting \\
Explanation & Transparent design + precedents & LIME / SHAP only & Faithful by construction \\
\bottomrule
\end{tabular}
\end{table}

\section{Chapter Conclusions}
\label{sec:sota_summary}

None of these building blocks is new. What is scarce, and what this work builds, is their disciplined
integration into a single explainable system for the retail investor: a transparent detector, semantic
retrieval paired with event-study impact, and explanations that are faithful by construction. The next
chapter turns these choices into a method.
